% Part: lambda-calculus
% Chapter: lambda-definability
% Section: partial-recursive-functions

\documentclass[../../../include/open-logic-section]{subfiles}

\begin{document}

\olfileid{lam}{ldf}{par}
\olsection{আংশিক পুনরাবৃত্ত অপেক্ষকগুলি \usetoken{S}{lambda definable}}

মৌলিক অপেক্ষকগুলি থেকে মিশ্রণ, আদিম পুনরাবৃত্তি ও সীমাহীন ন্যূনীকরণের
সাহায্যে যেসব অপেক্ষক পাওয়া যায়, সেগুলিই আংশিক পুনরাবৃত্ত অপেক্ষক।
সাধারণ পুনরাবৃত্ত অপেক্ষক থেকে এদের পার্থক্য হলো, সীমাহীন অনুসন্ধানে
ব্যবহৃত অপেক্ষকগুলিকে নিয়মিত হতে হয় না। নিয়মিততা আবশ্যক না হওয়ায়
ন্যূনীকরণ দিয়ে সংজ্ঞায়িত অপেক্ষক কখনও কখনও অসংজ্ঞায়িত হতে পারে।

প্রথমে মনে হতে পারে, (সর্বত্র-সংজ্ঞায়িত) সাধারণ পুনরাবৃত্ত অপেক্ষকগুলি
সবই !!{lambda definable} দেখাতে ব্যবহৃত একই পদ্ধতিতে প্রমাণ করা যাবে যে
সব আংশিক পুনরাবৃত্ত অপেক্ষকও !!{lambda definable}। যেমন, $f$ ও~$g$-কে
যথাক্রমে $F$ ও~$G$ পদ !!{lambda define} করলে, $g$-এর সঙ্গে $f$-এর
মিশ্রণকে $\lambd[x][F (G x)]$ !!{lambda define} করে। কিন্তু অপেক্ষকগুলি
আংশিক হলে এতে সমস্যা হয়। $g(x)$ অসংজ্ঞায়িত হলে $G x$-এর কোনো স্বাভাবিক
রূপ থাকে না। অধিকাংশ ক্ষেত্রে এর অর্থ $F (G x)$-এরও কোনো স্বাভাবিক রূপ
নেই, আর সেটিই আমরা চাই। কিন্তু $F$ যদি $\lambd[x][\lambd[y][y]]$ হয়,
তাহলে $F (G x)$-এর একটি স্বাভাবিক রূপ ($\lambd[y][y]$) আছে।

এই সমস্যা অতিক্রম করা অসম্ভব নয়; সব আংশিক পুনরাবৃত্ত অপেক্ষককে এমনভাবে
!!{lambda define} করার উপায় আছে, যাতে অসংজ্ঞায়িত মানগুলি স্বাভাবিক রূপহীন
পদ দিয়ে উপস্থাপিত হয়। কিন্তু সাধারণ পুনরাবৃত্ত অপেক্ষকের জন্য নেওয়া
পদ্ধতির চেয়ে এই উপায়গুলি কিছুটা বেশি জটিল ও কম স্বজ্ঞামূলক। এখানে
প্রমাণ ছাড়াই উপপাদ্যটি লিপিবদ্ধ করি:

\begin{thm}
  সব আংশিক পুনরাবৃত্ত অপেক্ষক !!{lambda definable}।
\end{thm}

\end{document}
